\begin{abstract}
    On-policy distillation (OPD) combines student-generated rollouts with dense token-level supervision from a teacher. Existing work has mainly studied its algorithmic behavior, leaving the role of training data unclear. We examine this role at the data-minimal limit by training on a single query. One-shot OPD keeps improving for hundreds of steps and recovers most of full-data OPD's gain across task domains and model families. We explain this result through the states visited during training and the rate at which the student aligns with the teacher. We measure \emph{state coverage}, the fraction of the states full-data OPD visits that a query set's rollouts reach. A single query already reaches \(71.5\%\), most of it within the first 100 steps. Adding semantically distinct queries raises coverage and validation accuracy together, until 16 queries reach \(98.9\%\) and match full-data training. Yet alignment slows at a similar pace whether OPD trains on one query or the whole dataset, and even a fixed set of states takes hundreds of steps to absorb. OPD is therefore data-overfed but algorithm-starved. Its rollouts quickly expose broad supervision, while the student absorbs that supervision increasingly slowly. The state-coverage result extends to multi-teacher OPD, where 16 semantically diverse queries per domain match full-data MOPD. As a further stress test, content-light templates and off-domain WildChat queries also approach the real-query baseline. Task content and induced state coverage can therefore come apart. We hope these findings direct future work toward the step efficiency of OPD, and prompt a re-examination of the data and the mechanisms behind its recent successes in frontier post-training.
% On-policy distillation (OPD) combines student-generated rollouts with dense token-level supervision from a teacher. Existing work has mainly studied its algorithmic behavior, leaving the role of training data unclear. We examine this role at the data-minimal limit by training on a single query. One-shot OPD keeps improving for hundreds of steps and recovers most of full-data OPD's gain across task domains and model families. We explain this result through the states visited during training and the rate at which the student aligns with the teacher. A single query reaches 71.5\% state coverage relative to full-data OPD, with most coverage appearing in the first 100 steps. Adding semantically distinct queries increases state coverage and validation accuracy together. 16 queries reach 98.9\% coverage and match full-data training. Yet alignment slows at a similar pace whether OPD trains on one query or the whole dataset, and even a fixed set of states takes hundreds of steps to absorb. \textbf{OPD is therefore data-overfed but algorithm-starved.} Its rollouts quickly expose broad supervision, while the student absorbs that supervision increasingly slowly. The state-coverage result extends to multi-teacher OPD, where 16 semantically diverse queries per domain match full-data MOPD. As a further stress test, content-light templates and off-domain WildChat queries also approach the real-query baseline. These findings show that task content and induced state coverage can partly come apart. These findings can inspire future work on the step efficiency of OPD, and encourage a re-examination of the data and the mechanisms behind its recent successes in frontier post-training.
\end{abstract}

% On-policy distillation (OPD) trains a student on its own rollouts while a teacher returns a full next-token distribution at every visited prefix, yet existing accounts of OPD's efficiency are algorithm-side, leaving unstudied how the training data shapes OPD and how the two interact.
% We therefore take OPD to its data-minimal limit and train on a single query (one-shot OPD).
% One-shot OPD keeps improving for hundreds of steps and, over a thousand steps of training, recovers 72\% of full-data OPD's gain in mathematics; the effect holds across task domains and model families, and persists even when the student never solves the training query.
% This poses the question the paper answers: \textit{Why can OPD, trained on just one example, keep learning for so long and improve by so much?}
% Analyzing it from the data and the algorithm sides gives one answer: \textbf{OPD is data-overfed but algorithm-starved}, and one query already supplies more supervision than the algorithm can absorb.
% On the data side, \textbf{what OPD consumes is states, not queries}: one query's rollouts reach 71.5\% of the state space that training on DAPO-Math-17K explores, most of it within the first 100 steps, and semantically distinct queries raise state coverage and validation accuracy together only until 16 of them reach 98.9\% state coverage and match full-data OPD.
% On the algorithm side, \textbf{what paces a run is absorption, not supply}: the student closes the remaining teacher--student gap ever more slowly, absorbing four to six times less of what is left per update at step 200 than at step 50, and it slows at the same pace from one query to all 17k, so the pace of a run is a property of OPD rather than of the training set.
% The same query-diversity result carries over to multi-teacher OPD, where one student trains on several domains at once, each with its own teacher: 16 semantically diverse queries per domain match full-data MOPD.
% Taken further, the training input need only lead the student into reasoning, since off-domain WildChat queries and even content-free templates also drive effective OPD.
% \textbf{What data design has to settle is therefore no longer how many queries are collected, but which states they induce for capable teachers to supervise.}
% These findings can inspire future work on the step efficiency of OPD, and encourage a re-examination of the data and the mechanisms behind its recent successes in frontier post-training.
